% ตัวอย่าง มคอ.5 — รายงานการสัมฤทธิ์ผลของผู้เรียนระดับรายวิชา
% Compile: xelatex example.tex (สองครั้งถ้ามี cross-ref)
\documentclass{mko5}

\begin{document}

\mkofivetitle

% ----- ข้อมูลหัวรายงาน -----
\coursecode{7015101}
\coursename{เทคโนโลยีอุบัติใหม่ทางวิศวกรรมคอมพิวเตอร์และปัญญาประดิษฐ์}
\curriculumnote{หลักสูตรวิศวกรรมศาสตรมหาบัณฑิต สาขาวิชาวิศวกรรมคอมพิวเตอร์และปัญญาประดิษฐ์ พ.ศ. 2568}
\coursegroup{วิชาเฉพาะ}
\coursecredits{3 (3-0-6)}
\studyplan{แผน 2 (วิชาชีพ --- สารนิพนธ์)}
\semester{ภาคเรียนที่ 2 / ปีการศึกษา 2568}
\prerequisite{ไม่มี}
\corequisite{ไม่มี}
\studyplace{มหาวิทยาลัยราชภัฏอุตรดิตถ์ (คณะเทคโนโลยีอุตสาหกรรม / ออนไลน์ผสม)}
\registeredstudents{4 คน}
\completedstudents{4 คน}
\instructor{ผศ.ดร.พิศิษฐ์ นาคใจ, ผศ.ดร.พัชรี มณีรัตน์}
\coinstructors{---}

\mkoheaderinfo

% ----- 1. กิจกรรมที่ดำเนินการจริง -----
\begin{teachingactualtable}
  \actualrow{CLO1}{บรรยายสาระสำคัญเทคโนโลยีอุบัติใหม่ (AI, IoT, Blockchain) พร้อมกรณีศึกษา;
  อภิปรายเชิงวิพากษ์; มอบหมายให้อ่านบทความวิจัยล่าสุด}
  \actualrow{CLO2}{Workshop การประยุกต์ใช้เครื่องมือ ML/DL บนชุดข้อมูลจริง;
  ปฏิบัติการกลุ่มย่อย; นำเสนอผลงานต้นแบบ}
  \actualrow{CLO3}{สัมมนาเชิงปฏิบัติการ IEL Community Workshop (ถ่ายทอดความรู้ AI ให้โรงเรียนใน จ.อุตรดิตถ์)}
\end{teachingactualtable}

% ----- 2. วิธีการวัดผลที่ใช้ -----
\begin{assessmentusedtable}
  \assessusedrow{CLO1}{แบบทดสอบกลางภาค, รายงานกลุ่ม, Rubric 4 ระดับ}
  \assessusedrow{CLO2}{โครงงานปฏิบัติ, การนำเสนอ, แบบประเมินตนเองและเพื่อน}
  \assessusedrow{CLO3}{รายงานกิจกรรมบริการวิชาการ, สังเกตพฤติกรรม, แบบสะท้อนคิด}
\end{assessmentusedtable}

% ----- 3. สรุปการสัมฤทธิ์ผล -----
\mkosectionthree

\closummaryblock{CLO1}{อธิบายและวิเคราะห์เทคโนโลยีอุบัติใหม่ที่เกี่ยวข้องกับวิศวกรรมคอมพิวเตอร์ได้อย่างถูกต้อง}{80}
\achievementpass{4}{100}{ผ่านเกณฑ์ระดับ 3 ขึ้นไปทุกคน}
\achievementfail{0}{0}{ไม่มี}
\closummaryend

\closummaryblock{CLO2}{ออกแบบและพัฒนาต้นแบบระบบจากโจทย์ที่กำหนดได้}{80}
\achievementpass{4}{100}{ผ่านการประเมินโครงงานและการนำเสนอ}
\achievementfail{0}{0}{ไม่มี}
\closummaryend

\closummaryblock{CLO3}{ถ่ายทอนความรู้ด้านปัญญาประดิษฐ์ให้ชุมชน/โรงเรียนได้อย่างมีประสิทธิภาพ}{80}
\achievementpass{4}{100}{เข้าร่วม IEL Workshop ครบและผ่าน Rubric}
\achievementfail{0}{0}{ไม่มี}
\closummaryend

% ----- 4. ปัญหา/อุปสรรค -----
\mkosectionfour
\begin{mkotextblock}
จำนวนนักศึกษาน้อย ($n=4$) ทำให้การวิเคราะห์สถิติเชิงเปรียบเทียบมีข้อจำกัด;
บางสัปดาห์ต้องจัดการเรียนแบบออนไลน์เนื่องจากภารกิจบริการวิชาการนอกสถานที่
\end{mkotextblock}

% ----- 5. แผนปรับปรุง -----
\mkosectionfive
\begin{mkotextblock}
เพิ่มชุดข้อมูลฝึกปฏิบัติที่หลากหลายขึ้น; จัดทำแบบฟอร์มสะท้อนคิดมาตรฐานสำหรับ CLO3;
ประสานกับหลักสูตรเพื่อรวบรวมผล มคอ.5 ทุกวิชาเข้าสู่การทบทวนระดับหลักสูตรภาคถัดไป
\end{mkotextblock}

% ----- 6. ข้อเสนอแนะ -----
\mkosectionsix
\begin{mkotextblock}
ขอให้คณะสนับสนุนการเข้าถึง GPU/ห้องปฏิบัติการสำหรับรายวิชา CPE\&AI;
พิจารณาจัดอบรม Rubric-based assessment สำหรับอาจารย์ผู้สอนรายวิชาใหม่
\end{mkotextblock}

% ----- 7. ลงลายมือชื่อ -----
\mkosectionseven
\mkosignaturepair{ผศ.ดร.พิศิษฐ์ นาคใจ}{27 มีนาคม 2569}{ผศ.ดร.พิทักษ์ คล้ายชม}{27 มีนาคม 2569}

\end{document}
